% !PRINTED 196
\weberpage{196}
%<<CONT
and we regard our problem as solved when a representative of each
of these families has been determined.

We shall later make use of this freedom in order to represent
the groups found as simply as possible.

The determinant
\[
\alpha\delta - \beta\gamma = \Delta
\]
must be different from zero, and if simplicity demands it,
we may indeed assume it $= 1$, which we shall occasionally
do.

In the sense of group theory we denote the repetition
of a substitution by exponents: $\Theta, \Theta^2, \Theta^3, \ldots$,
by which, then, are to be understood not powers, but always
linear substitutions of the group (1) again. The identical
substitution $x = x$, which is to be regarded as the unit of the
group, is also denoted by $\Theta^0$ or by 1.

Two inverse substitutions $\Theta, \Theta^{-1}$ can be represented
in the form:
\[
\Theta = \begin{pmatrix} \alpha, & \beta \\ \gamma, & \delta \end{pmatrix}, \quad
\Theta^{-1} = \begin{pmatrix} \delta, & -\beta \\ -\gamma, & \alpha \end{pmatrix}.
\]

Since the group is by hypothesis finite, each of its
elements has a definite \gesperrt{degree}, i.e.\ for each $\Theta$ there
is a definite least positive number $t$ for which $\Theta^t = 1$.
This number $t$ must be a divisor of $n$.

For what follows it is of importance to consider, for each
of the substitutions of the group (excepting the identical
substitution), the values of the variable that become equal to
their transforms, that is, the roots of the equations
\weq{5}{x = \Theta(x).}
These values we shall call the \gesperrt{poles} of the substitution $\Theta$.

The equation (5) is quadratic and, if one substitutes for
$\Theta$ the expression (2), takes the form:
\weq{6}{\gamma x^2 + (\delta - \alpha) x - \beta = 0.}
This equation has two roots, which are equal to one another
only when
\weq{7}{(\delta - \alpha)^2 + 4\beta\gamma = 0}
or
\weq{8}{(\alpha + \delta)^2 = 4\Delta}
%CONT>>
